\subsection{Definition of barycenters}

Let E be a Hausdorff locally convex space over R, \(E'\) its dual, and \(E'^{*}\) the algebraic dual of \(E'\) , E being canonically identified with a linear subspace of \(E'^{*}\) . Let K be a compact subset of E; the canonical injection of K into E being continuous with compact support, for every measure \(\mu\) on K the integral \(\int x d\mu(x)\) is therefore defined and is an element of \(E'^{*}\) (Ch. III, §3, No.~1). Moreover, on K, the topology induced by the weak topology \(\sigma(E'^{*}, E')\) is identical with the original topology. Finally, if C is the closed convex envelope of K in \(E'^{*}\) equipped with \(\sigma(E'^{*}, E')\) , then \(C \cap E\) is the closed convex envelope of K in E for the original topology (or for the weakened topology \(\sigma(E, E')\) ).

DEFINITION 1. --- Let K be a compact subset of a Hausdorff locally convex space E. For every positive measure \(\mu\) on K of total mass equal to 1, the vector \(\mathbf{b}_{\mu} = \int \mathbf{x} d\mu(\mathbf{x})\) (belonging to \(E'^{*}\) ) is called the barycenter of \(\mu\) .

Example. --- Let \(\mu\) be a discrete measure on K, positive and of total mass 1; it is thus of the form \(\mu = \sum_{i=1}^{n} \lambda_i \varepsilon_{x_i}\) , where \(x_i \in K\) , and the \(\lambda_i\) are real numbers such that \(\lambda_i \geqslant 0\) for all i and \(\sum_{i=1}^{n} \lambda_i = 1\) . Since \(\int x d\varepsilon_y(x) = y\) (Ch. III, §3, No.~1, Example 3), \(b_\mu = \int x d\mu(x) = \sum_{i} \lambda_i x_i\) . In particular, for every \(x \in K\) , x is the barycenter of the measure \(\varepsilon_x\) .

PROPOSITION 1. --- Let E be a Hausdorff locally convex space, K a compact subset of E, and C the closed convex envelope of K in E. The set C then consists of the points of E that are barycenters of at least one positive measure of mass 1 on K.

This is nothing more than Prop. 5 of Ch. III, §3, No.~2 applied to the canonical injection of K into E.

COROLLARY. --- If the closed convex envelope C of K in E is compact, then the barycenter of every positive measure of total mass 1 on K belongs to E.

For, C is then also the closed convex envelope of K in E'* equipped with the weak topology \(\sigma (\mathrm{E}^{\prime *},\mathrm{E}^{\prime})\) , and it suffices to apply, to the canonical injection of K into E, Prop. 4 of Ch. III, §3, No.~2.

Remark. --- The Cor. of Prop. 1 is applicable in particular when K is convex or when E is quasi-complete.

PROPOSITION 2. --- Let K be a compact convex subset of a Hausdorff locally convex space E, \(\mu\) a positive measure of total mass 1 on K, \(b_{\mu}\) its barycenter. For every convex numerical function \(f \geqslant 0\) that is lower semicontinuous on K,

\[
f (\mathbf {b} _ {\mu}) \leqslant \int^ {*} f d \mu .
\]

It is known (TVS, II, §5, No.~4, Prop. 5) that \(f\) is the upper envelope of a family of restrictions to \(K\) of continuous affine linear functions \(h_{\alpha} : \mathbf{x} \mapsto c_{\alpha} + \langle \mathbf{x}, \mathbf{z}_{\alpha}' \rangle\) . Then

\[
\int h _ {\alpha} (\mathbf {x}) d \mu (\mathbf {x}) \leqslant \int^ {*} f (\mathbf {x}) d \mu (\mathbf {x})
\]

for every \(\alpha\) ; now, \(\int h_{\alpha}(\mathbf{x}) d\mu (\mathbf{x}) = c_{\alpha} + \int \langle \mathbf{x}, \mathbf{z}_{\alpha}' \rangle d\mu (\mathbf{x})\) since \(\mu\) has total mass 1; but \(\int \langle \mathbf{x}, \mathbf{z}_{\alpha}' \rangle d\mu (\mathbf{x}) = \langle \mathbf{b}_{\mu}, \mathbf{z}_{\alpha}' \rangle\) by the definition of barycenter. Therefore \(\sup \int h_{\alpha}(\mathbf{x}) d\mu (\mathbf{x}) = \sup h_{\alpha}(\mathbf{b}_{\mu}) = f(\mathbf{b}_{\mu})\) , whence the conclusion.

When \(\mu\) is a discrete positive measure on K of total mass 1, Prop. 2 yields anew the inequality that defines the convex functions on K.

COROLLARY. --- For every convex numerical function \(g\) on \(K\) that is bounded and lower semi-continuous, \(g(\mathbf{b}_{\mu}) \leqslant \int g d\mu\) .

It suffices to note that \(\inf_{\mathbf{x} \in \mathrm{K}} g(\mathbf{x}) = a\) is finite and apply Prop. 2 to \(g - a\) .
% semantic-fragments: legacy-input-seam
\relax{}
